\section{The precise constructive claim}
\label{sec:foundations}

The setting is first-order intuitionistic primitive recursive arithmetic,
with a symbol for each primitive recursive definition, its defining
equations, successor separation and injectivity, intuitionistic predicate
logic with equality, and induction for quantifier-free formulas. This is
the system \code{PRAAx} in the inspected \code{ConciseIPRA/PR.lean}.
Its arithmetic interpretation in the native concise calculus is supplied by
\code{pra\_iff\_native} in \code{NativeBiInterpretation.lean}~\cite{repo}.
We work in this ordinary presentation to make the mathematics readable.

There is no sort of arbitrary real numbers or arbitrary real functions in
this theory. Accordingly, an assertion about a real or complex argument
means a theorem schema for \emph{each fixed primitive recursive name}, with
the displayed numerical domain information and error certificates. The
rational-argument estimates are uniform first-order statements about
rational codes and precisions. Substituting a fixed name and its modulus
gives the corresponding theorem about that name. We do not postulate a
primitive recursive universal evaluator for all primitive recursive
programs; no such evaluator exists.

\subsection{Finite arithmetic and names}

Integers are sign-and-magnitude codes, and a rational is a pair $(u,v)$
with $u\in\Z$ and $v\in\N_{>0}$. Equality and order are cross-multiplication
tests. Reduction by gcd is convenient but unnecessary. Finite lists,
factorials, powers, sums, products, finite differences, and bounded search
are primitive recursive. Complex rationals are pairs of rationals. For
coding we use $\norm{x+iy}_1=|x|+|y|$; estimates written with the ordinary
complex modulus can instead be checked after squaring, or transferred with
$|z|\le\norm z_1\le2|z|$. Enlarge a disk radius by a factor of two when
using the latter convention.

\begin{definition}[Names and equality]
A regular real name is a fixed primitive recursive map $a:\N\to\Q$
satisfying
\[
 |a(p)-a(q)|\le2^{-p}+2^{-q}\qquad(p,q\in\N).
\]
A complex name is a pair of such names; alternatively use this inequality
in the $1$-norm. Equality of names is
\[
 a\simeq b\quad\Longleftrightarrow\quad
 \forall p\; |a(p+2)-b(p+2)|\le2^{-p}.
\]
An inequality $a\le b$ means the analogous comparison with error
$2^{-p}$. Positive separation from zero is supplied as a rational number
$d>0$ and the certificate $a\ge d$. Complex apartness is supplied as
$|a|\ge d$; it is not inferred from $\neg(a=0)$.
\end{definition}

These equalities are equivalence relations. For example, to prove
transitivity at precision $p$, use the hypotheses at precision $p+r$,
regularity to compare stages, and take $r\ge4$; to remove a residual
positive rational error, assume the desired rational inequality fails and
choose a precision bounded by the binary length of its positive gap.
This is a bounded rational argument, not excluded middle for equality of
names. Equivalently, one may define equality by
$\forall p,q\;|a(p)-b(q)|\le2^{-p}+2^{-q}$; the same refinement argument
proves the equivalence of the conventions.

A uniformly convergent finite expression $F_N$ is represented by
\emph{substitution of a modulus}: if
$|F_M-F_N|\le C/N$ for $M\ge N\ge N_0$, use
$N(p)=N_0+\ceil C\,2^{p+4}$ and evaluate $F_{N(p)}$ to error
$2^{-p-4}$. The output is a regular name after a fixed precision shift.
This defines an actual primitive recursive function. It does not invoke
a principle saying that an arbitrary Cauchy sequence has a limit.

\subsection{How the displayed proofs become arithmetic proofs}

Every finite sum identity is proved by induction on its finite length.
Every denominator has a displayed positive lower bound. A bounded
quantifier over a decidable predicate is represented by a finite Boolean
product or sum, hence by a primitive recursive characteristic function.
Thus the invariant to which induction is applied can be an equation
between primitive recursive terms. When several variables or derivative
orders are present, they remain parameters of this quantifier-free
induction.

There is one essential discipline. We never use induction on a statement
such as ``two names are equal,'' which already contains an unbounded
precision quantifier. We first fix the requested precision, replace every
name in the finite calculation by a rational approximation, and allocate
its error budget using the following elementary estimates:
\begin{align}
 \left|\sum_{j<s}a_j-\sum_{j<s}b_j\right|&\le s\varepsilon,
 \label{eq:sum-budget}\\
 \left|\prod_{j<s}a_j-\prod_{j<s}b_j\right|
 &\le sM^{s-1}\varepsilon,
 \quad |a_j|,|b_j|\le M,
 \label{eq:product-budget}\\
 |a^{-1}-b^{-1}|&\le\varepsilon/d^2,
 \quad |a|,|b|\ge d.\label{eq:inverse-budget}
\end{align}
Here $|a_j-b_j|\le\varepsilon$. The product estimate follows by replacing
one factor at a time, so its proof is a finite telescoping identity. When
only $|a|\ge d$ is known, evaluate $b$ with error at most $d/2$ and use
$2\varepsilon/d^2$ instead. Integer upper bounds for all constants suffice.
For an integer $C\ge1$, the unnecessarily large shift $C+k+3$ makes
$C2^{-(C+k+3)}<2^{-k}$, eliminating any need for logarithmic search.

\begin{proposition}[Finite-precision reading]
The proofs below give, for each fixed tuple of names and each requested
precision, finite rational identities or inequalities with primitive
recursive lengths, cutoffs, and precision shifts. Their induction
invariants are quantifier-free. Consequently their displayed equality,
inequality, derivative, and integral conclusions have derivations in
$\ipra$ under the stated certificates.
\end{proposition}

\begin{proof}
Expand rational operations into integer operations, replace finite sums
and products by their recursors, and use
\eqref{eq:sum-budget}--\eqref{eq:inverse-budget} after each arithmetic
operation. All searches in the constructions have an explicit finite
length. The only case distinctions are on integers or rational
comparisons. A series tail is always replaced by a displayed bound before
the target precision is chosen. An integral is always replaced by a
finite sum before invoking a calculus identity. The specific invariants
and moduli are collected in \cref{sec:formalization}; the estimates are
proved in their mathematical sections. Finally discharge the precision
parameter by intuitionistic universal introduction. This is a proof
translation recipe for these concrete arguments, not a claim that
primitive recursiveness alone implies provability of arbitrary statements.
\end{proof}

\subsection{Elementary functions with explicit algorithms}

For $|w|\le B$, $B\ge1$ an integer, define $\exp w$ by the finite Taylor
polynomials. If $m+1\ge2B$, the remaining terms have successive absolute
ratios at most $1/2$, and therefore
\begin{equation}
 \left|\exp w-\sum_{j=0}^{m}\frac{w^j}{j!}\right|
 \le\frac{2B^{m+1}}{(m+1)!}.\label{eq:exp-tail}
\end{equation}
For example $m=4B+p+6$ makes this less than $2^{-p}$: the term at
$2B$ is at most $3^B$, since finite multinomial expansion and
$\sum_{j\ge0}1/j!<3$ give $\sum B^j/j!\le3^B$, and the subsequent
ratios are at most $1/2$. The bound $\sum1/j!<3$ itself follows from
$j!\ge2^{j-1}$ for $j\ge1$. All these are finite inequalities with
geometric remainder bounds. Finite binomial convolution proves
$\exp(u+v)=\exp u\exp v$. In particular $\exp u\ne0$ with explicit
lower bound $3^{-B}$ for $|u|\le B$.

For $q>0$, put $v=(q-1)/(q+1)$. Define
\begin{equation}
 \log q=2\sum_{j\ge0}\frac{v^{2j+1}}{2j+1},\qquad
 |\text{tail after }j<L|\le
 \frac{2|v|^{2L+1}}{(2L+1)(1-v^2)}.\label{eq:log-alg}
\end{equation}
If $B^{-1}\le q\le B$, then $|v|\le1-1/B$ and
$(1-1/B)^B\le1/2$ (use the finite binomial inequality on the
reciprocal). Thus $L=B(p+2B+4)$ is a safe primitive recursive cutoff.
Finite differentiation of these series gives $(\log q)'=1/q$,
$(\exp q)'=\exp q$; the derivative tails are geometric on slightly
smaller intervals. The product and inverse identities, and
$\log(ab)=\log a+\log b$ for positive $a,b$, follow either by finite
series convolution or by integrating their zero derivative on a compact
interval. Rational or complex powers of a \emph{positive real} $t$ mean
$t^z=\exp(z\log t)$. This convention has no complex branch ambiguity.

Define sine and cosine by their Taylor series, equivalently by
$(\exp(iw)-\exp(-iw))/(2i)$ and
$(\exp(iw)+\exp(-iw))/2$. The addition identities follow from the
exponential identity. Define
$\pi=4\int_0^1(1+t^2)^{-1}\dd t$; the alternating arctangent sum has
error at most $4/(2L+1)$, so $L=2^{p+4}$ suffices.
For completeness the identification with trigonometric period is finite:
write $a(u)=\int_0^u(1+t^2)^{-1}\dd t$, differentiate and check
\[
 \sin(2a(u))=\frac{2u}{1+u^2},\qquad
 \cos(2a(u))=\frac{1-u^2}{1+u^2}.
\]
Both pairs have the same initial values and satisfy the same two first
order equations. Uniqueness here requires no general differential
equation theorem: the squared difference has zero derivative, and the
finite fundamental theorem described below makes it zero. At $u=1$ this
gives $\sin(\pi/2)=1$, $\cos(\pi/2)=0$ and then
$\sin\pi=0$, $\cos\pi=-1$. The map $u\mapsto2a(u)$ has derivative
at least one on $[0,1]$. The approximate-sign rational bisection described
in the trigonometric section therefore constructs its inverse on
$[0,\pi/2]$ with a prescribed rate. Positivity on the first quadrant
follows from this parametrization. Concavity, $\sin''=-\sin\le0$, gives
\begin{equation}
 \sin v\ge2v/\pi\quad(0\le v\le\pi/2).\label{eq:sine-chord}
\end{equation}
These facts, addition and periodicity give the needed finite lists of
trigonometric roots. The coarse rational bounds $2<\pi<4$ follow directly
from the defining integral.

\subsection{The finite calculus convention}
\label{sec:finite-calculus}

For a function on a rational interval $[a,b]$ with a supplied Lipschitz
constant $L$, the left sum on $m$ equal subintervals has error at most
$L(b-a)^2/m$ against any refinement. Use
$m\ge2^{p+4}(1+\ceil{L(b-a)^2})$ and evaluate each summand with
error $2^{-p-4}/(1+b-a)$. The resulting diagonal sequence defines
$\int_a^b f$ as a name.

The finite fundamental theorem follows by telescoping
$F(t_{j+1})-F(t_j)$. If a derivative approximation satisfies
\[
 |F(t+h)-F(t)-h f(t)|\le C h^2,
\]
the sum of its errors is at most $C(b-a)^2/m$. This proves
$\int_a^b f=F(b)-F(a)$ when such a quadratic remainder is supplied.
It applies to our elementary expressions by the finite Taylor series
bounds, and to the gamma expressions by the derivative estimates below.
The product rule with its quadratic remainder proves integration by
parts. Applying this telescoping argument to $F(\phi(t))$, with
explicit bounds for $\phi'$ and its remainder, proves substitution.

Multiple integrals over rectangles are iterated finite sums; changing
their order is a finite rearrangement. We do not need a general
multidimensional change-of-variables theorem: the beta--gamma identity
below is proved by a one-dimensional shift and a finite-product limit.
Singular edges and infinite tails are cut off first using the estimates
in \cref{sec:integrals}. Thus no use is made of Fubini or dominated
convergence for arbitrary functions.

Derivatives of order $r$ mean local polynomial expansions with a
primitive recursive remainder bound. Holomorphic means the corresponding
complex power expansions on every certified disk; entire means that
these bounds are available for every integer radius. Meromorphic means
that on a certified neighborhood a fixed finite power of $(z-z_0)$ times
the function has such an analytic presentation. These are constructive
data, not tests for deciding analyticity of an arbitrary program.
